% 共通マクロ集（レポートテンプレート用）
% 必要に応じてプロジェクトごとに上書き・拡張してください。

% ===== siunitx の設定 =====
% 単位・数値表示のための設定（siunitx は main.tex 側で読み込まれている前提）
\sisetup{detect-all=true, separate-uncertainty=true}

% ===== 参照用ショートカット =====
% 図・表・式を参照する際に使用
% 使用例: \figref{fig:example} → 図\ref{fig:example}
\newcommand{\figref}[1]{図\ref{#1}}
\newcommand{\tabref}[1]{表\ref{#1}}
\newcommand{\eqnref}[1]{式(\ref{#1})}

% 添字をローマン体にするマクロ
\newcommand{\subr}[1]{_{\mathrm{#1}}}

% ===== 注記用マクロ =====
\newcommand{\mynote}[1]{\textbf{[注]}\,\textit{#1}}

% ---- 訂正表示 ----
\newcommand{\cor}[1]{\textcolor{red}{#1}}
% ===== 図のキャプションや浮動体の共通スタイル =====
% 必要に応じてキャプションのフォーマットを設定
% 例: \usepackage{caption}
%     \captionsetup{font=small, labelfont=bf}

% ===== フロントマター用マクロ =====
% 表紙に埋め込む情報を簡単に更新するためのマクロ
\newcommand{\ReportTitle}{レポートタイトルをここに記入}
\newcommand{\ReportSubtitle}{実験テーマ・解析対象・副題など}
\newcommand{\ReportCourse}{講義名／実験名}
\newcommand{\ReportDepartment}{化学物理工学科}
\newcommand{\ReportClass}{3年}
\newcommand{\ReportStudentID}{24264027}
\newcommand{\ReportAuthor}{久保木 琢真}
\newcommand{\ReportCollaborators}{共同実験者なし}
\newcommand{\ReportSupervisor}{担当教員名}
\newcommand{\ReportInstitution}{東京農工大学 工学部 化学物理工学科}
\newcommand{\ReportDate}{\today}

% ===== 分野別パッケージ =====
% 必要に応じて以下のコメントを外してパッケージを読み込む

% 化学式が必要な場合:
% \usepackage[version=4]{mhchem}
% 使用例: \ce{H2O}, \ce{CO2}

% 回路図が必要な場合:
% \usepackage{circuitikz}

% 化学構造式が必要な場合:
% \usepackage{chemfig}

% ===== カスタムコマンド例 =====
% プロジェクトに応じて独自のコマンドを定義できます
% 例: \newcommand{\diff}{\mathrm{d}}  % 微分記号
%     \newcommand{\subA}{実験A}  % サブセクション名の統一
